\chapter{European and Japanese Government Bonds}\label{ch:m2:european-and-japanese-government-bonds}

On Wednesday 28 September 2022 the yield of the thirty-year UK government bond
moved within a range of 127 basis points in a single day, more than its range
over a whole year in all but four of the previous twenty-seven. In the days
before, it had risen by 160 basis points, having started the year near 1.2\%.
Pension funds that had borrowed against their \omterm{def:m2:european-and-japanese-government-bonds:ldi}{gilts} to match their long
liabilities were being asked for collateral faster than they could raise it,
and their managers were telling the Bank of England that, at the prevailing
yields, some of their funds would be worth less than nothing by the next
morning. That day the Bank announced temporary purchases of long \omterm{def:m2:european-and-japanese-government-bonds:ldi}{gilts} to
restore orderly market conditions; by 14 October it had bought 19.3 billion
pounds' worth. The euro area, Japan and the United Kingdom each run their
government bond markets differently, and each has, in the last fifteen
years, had to decide what its central bank owes the market in a crisis. This
chapter compares them.

\section{The euro area: one currency, many issuers}

Twenty-one governments borrow in euros since Bulgaria joined on 1 January
2026, and none of them controls the central bank that issues it. The same currency therefore trades at many yields, and the
differences are the market's measure of each government's credit and of the
risk that the currency union itself might change.

\begin{definition}[Sovereign spread and redenomination risk]\label{def:m2:european-and-japanese-government-bonds:spread}
A \emph{sovereign spread}\index{sovereign spread} is the difference between the
yield of one government's bond and that of a benchmark government of the same
currency and maturity, in the euro area usually Germany. \emph{Redenomination
risk}\index{redenomination risk} is the part of a spread that pays for the
possibility that a bond is repaid in a new national currency worth less than
the euro.
\end{definition}

\begin{definition}[Bund and syndication]\label{def:m2:european-and-japanese-government-bonds:bund}
A \emph{Bund}\index{Bund} is a German federal government bond (Bundesanleihe),
issued at 7, 10, 15 or 30 years with an annual coupon; the ten-year Bund's
yield is the euro area's benchmark long rate. \emph{Syndication}\index{syndication}
is the sale of a new bond through a group of banks that build a book of
investor orders and set the price, instead of an auction; some issuers use it
for new long bonds, where demand is hardest to predict.
\end{definition}

\begin{omfigure}
\begin{tikzpicture}
\begin{axis}[width=11cm,height=5cm,
  ylabel={spread to Germany (bp)},
  tick label style={font=\small},label style={font=\small},
  xmin=2007,xmax=2027,ymin=-20,ymax=560,grid=major,grid style={gray!25},
  xtick={2008,2011,2014,2017,2020,2023,2026},
  xticklabel style={/pgf/number format/1000 sep=},
  legend columns=3,legend style={font=\footnotesize,at={(0.5,-0.18)},anchor=north,draw=none,fill=none,/tikz/every even column/.append style={column sep=8pt}},
  table/col sep=comma]
\addplot[omThm,very thick] table[x=t,y=IT]{figdata/markets-2/07-european-and-japanese-government-bonds/spreads.csv};
\addplot[omProp,thick] table[x=t,y=ES]{figdata/markets-2/07-european-and-japanese-government-bonds/spreads.csv};
\addplot[omDef,thick,dashed] table[x=t,y=FR]{figdata/markets-2/07-european-and-japanese-government-bonds/spreads.csv};
\legend{Italy,Spain,France}
\node[font=\footnotesize,anchor=west] at (axis cs:2014.3,480) {Italy, Nov 2011: 518 bp};
\end{axis}
\end{tikzpicture}

\omcaption{Ten-year spreads to Germany, monthly averages, 2007 to August 2026.
The sovereign debt crisis took Italy's to 518 basis points in November 2011
and Spain's to 555 in July 2012; spreads fell after the ECB announced outright monetary transactions in
September 2012. By August 2026 France's spread, 82 basis points, was slightly wider than
Italy's, 81. Data: OECD, via FRED.}\label{fig:m2:european-and-japanese-government-bonds:spreads}
\end{omfigure}

\begin{example}[Anatomy of a widening]\label{ex:m2:european-and-japanese-government-bonds:2011}
Between June and November 2011 (monthly averages) Italy's ten-year spread to
Germany widened by 326 basis points. Italy's yield rose 224 basis points; the
other 102 came from Germany's yield \emph{falling} by as much, as investors fled
into \omterm{def:m2:european-and-japanese-government-bonds:bund}{Bunds}. A spread measures relative credit, but half of a crisis widening
can be the benchmark rallying: a hedge that shorts the benchmark against the
spread loses on that half.
\end{example}

\section{Spreads, backstops and the electronic dealer market}

\begin{definition}[Central-bank backstop]\label{def:m2:european-and-japanese-government-bonds:backstop}
A \emph{central-bank backstop}\index{central-bank backstop} is a commitment, or
a standing instrument, by which a central bank will buy a government's bonds
in a market dysfunction, so that prices reflect fundamentals rather than a
run. Its conditions decide what it can be used for: the conditions for
activating it, the size, and whether it is announced as unlimited.
\end{definition}

The ECB has two. Outright monetary transactions, announced in September 2012,
have no ex ante limit but are conditional on the country having an adjustment
programme with the European Stability Mechanism: the promise is as large as
needed, and its price is a programme. The transmission protection
instrument of July 2022 is conditional instead on four criteria of fiscal
and macroeconomic soundness, and aims at ``unwarranted, disorderly market
dynamics that pose a serious threat to the transmission of monetary policy'',
purchases concentrating on maturities of one to ten years. A spread that
widens because of fundamentals is outside both; a spread that widens because
of a run is, in principle, inside. Deciding which is which is the ECB's call.

The secondary market is organised like the US one in two tiers. Among dealers,
European government bonds trade on electronic platforms, such as MTS's
inter-dealer markets, one for each country; clients trade with dealers by
request for quote on \omterm{def:m2:the-treasury-market:idb}{dealer-to-client platforms}, such as MTS's own BondVision
(\cref{ch:m2:how-bonds-trade}).

\begin{dated}{2026-09}{Ten-year yields, August 2026 averages}\label{dat:m2:european-and-japanese-government-bonds:yields}
Germany 3.18\%, Spain 3.63\%, Italy 3.99\%, France 4.00\%; Japan 2.94\%; United
Kingdom 4.99\%. Spreads to Germany: Spain 45, Italy 81, France 82 basis points.
\end{dated}

\section{Japan: yield-curve control and its exit}

\begin{definition}[Yield-curve control]\label{def:m2:european-and-japanese-government-bonds:ycc}
\emph{Yield-curve control}\index{yield-curve control} is a policy in which the
central bank sets a target for a long-term government bond yield and buys, in
whatever amount necessary, to hold it there, in addition to setting the short
rate.
\end{definition}

The Bank of Japan introduced it in September 2016, committing to buy
government bonds so that ten-year yields would stay ``more or less at the
current level (around zero percent)''. A yield target inverts the usual
relation between price and quantity: the central bank no longer chooses how
much it buys, the market does, by selling to it at the target. As yields rose
around the world, defending the target meant buying whatever was offered; on
20 December 2022 the Bank widened the band around it from a quarter to half a
percentage point and offered to buy ten-year bonds at 0.5\% every business day.
On 19 March 2024 it declared that \omterm{def:m2:european-and-japanese-government-bonds:ycc}{yield-curve control} and the
negative rate had ``fulfilled their roles'' and returned to steering the short
rate alone. The ten-year yield, near zero for most of the control period,
averaged 2.94\% in August 2026 (\cref{fig:m2:european-and-japanese-government-bonds:jgb}).

\begin{omfigure}
\begin{tikzpicture}
\begin{axis}[width=11cm,height=4.8cm,
  ylabel={10-year JGB yield (\%)},
  tick label style={font=\small},label style={font=\small},
  xmin=2007,xmax=2027,ymin=-0.5,ymax=3.3,grid=major,grid style={gray!25},
  xtick={2008,2011,2014,2017,2020,2023,2026},
  xticklabel style={/pgf/number format/1000 sep=},
  table/col sep=comma]
\fill[omProp!12] (axis cs:2016.72,-0.5) rectangle (axis cs:2024.21,3.3);
\addplot[omDef,very thick] table[x=t,y=JP]{figdata/markets-2/07-european-and-japanese-government-bonds/jgb.csv};
\node[font=\footnotesize,text=omProp,align=center] at (axis cs:2020.45,2.6) {yield-curve control\\Sep 2016 -- Mar 2024};
\end{axis}
\end{tikzpicture}

\omcaption{The ten-year Japanese government bond yield, monthly averages, 2007
to August 2026. Under \omterm{def:m2:european-and-japanese-government-bonds:ycc}{yield-curve control} (shaded) the Bank held it near its
target; after the exit it rose to the highest levels of the period. Data:
OECD, via FRED; dates of the policy from the Bank of Japan's
statements.}\label{fig:m2:european-and-japanese-government-bonds:jgb}
\end{omfigure}

\section{The United Kingdom and the 2022 gilt crisis}

\begin{definition}[Gilt and liability-driven investment]\label{def:m2:european-and-japanese-government-bonds:ldi}
A \emph{gilt}\index{gilt} is a UK government bond; conventional gilts pay
semiannual coupons, index-linked gilts pay coupons and principal indexed to
inflation. \emph{Liability-driven investment}\index{liability-driven
investment} (LDI) is a strategy in which a defined-benefit pension fund
hedges the interest-rate and inflation sensitivity of its liabilities with
long gilts and swaps, often with leverage from repo and derivatives, so that it
can hedge more of its liabilities than it holds in bonds while keeping other,
higher-returning assets.
\end{definition}

Leverage is what made the hedge fragile. A fund that holds \omterm{def:m2:european-and-japanese-government-bonds:ldi}{gilts} worth twice
its capital, borrowing the other half in repo, loses twice as fast when \omterm{def:m2:european-and-japanese-government-bonds:ldi}{gilts}
fall. Its manager keeps a cushion of capital against losses and asks the
pension fund for more when it runs low; for pooled funds, with many small
pension schemes behind them, that request takes days. In late September 2022
the yields moved in hours. The funds' only fast source of cash was their \omterm{def:m2:european-and-japanese-government-bonds:ldi}{gilts}:
selling them pushed yields higher, which called for more cash, which forced
more sales (\cref{fig:m2:european-and-japanese-government-bonds:loop}). On 26
September managers told the Bank that, at face value, they would have to sell
at least 50 billion pounds of long \omterm{def:m2:european-and-japanese-government-bonds:ldi}{gilts} in a short time, against average
trading of 12 billion a day in those maturities.

\begin{omfigure}
\begin{tikzpicture}[font=\small,
  box/.style={draw,thick,rounded corners=2pt,align=center,minimum height=0.9cm,minimum width=3.4cm,inner sep=4pt},
  arr/.style={-{Stealth[length=2.5mm]},thick}]
\node[box,draw=omThm,fill=omThm!8] (y) at (0,1.8) {long gilt yields rise};
\node[box,draw=omDef,fill=omDef!8] (c) at (5.6,1.8) {LDI cushions shrink};
\node[box,draw=omProp,fill=omProp!8] (m) at (5.6,-0.6) {collateral and margin calls\\on repo and swaps};
\node[box,draw=gray,fill=gray!8] (s) at (0,-0.6) {forced sales of long gilts\\into a thin market};
\draw[arr] (y) -- node[above=11pt] {prices fall} (c);
\draw[arr] (c) -- node[right=4pt] {leverage rises} (m);
\draw[arr] (m) -- node[below=17pt] {cash needed now} (s);
\draw[arr] (s) -- node[left=4pt] {more supply} (y);
\node[font=\footnotesize,text=omThm,align=center] at (2.8,0.6) {the loop the Bank\\broke on 28 September};
\end{tikzpicture}

\omcaption{The collateral spiral of September 2022. Pension-fund LDI strategies
borrowed against long \omterm{def:m2:european-and-japanese-government-bonds:ldi}{gilts}; rising yields eroded their cushions, calls for
collateral could be met fastest by selling \omterm{def:m2:european-and-japanese-government-bonds:ldi}{gilts}, and the selling raised
yields further. A buyer that could absorb the sales without needing a return,
the central bank, stopped the loop.}\label{fig:m2:european-and-japanese-government-bonds:loop}
\end{omfigure}

The Bank's purchases were targeted and temporary: long \omterm{def:m2:european-and-japanese-government-bonds:ldi}{gilts} only, from 28
September to 14 October, 12.1 billion of conventional and 7.2 billion of
index-linked \omterm{def:m2:european-and-japanese-government-bonds:ldi}{gilts}, sold back to the market between November 2022 and January
2023. The Financial Policy Committee then set a standard: LDI funds should be
able to withstand a rise of about 250 basis points in \omterm{def:m2:european-and-japanese-government-bonds:ldi}{gilt} yields without
forced sales, against the 100 basis points that earlier stress tests had
assumed.

\begin{proposition}[How far a levered bond fund can fall]\label{prop:m2:european-and-japanese-government-bonds:cushion}
A fund holds bonds worth $G$ with \omterm{def:m2:government-bonds:duration}{modified duration} $D$ and convexity $\mathcal C$,
financed by repo $B = \ell G$, so that its cushion is $(1-\ell)G$. Its cushion is
exhausted by the parallel yield rise $\Delta y^*$ at which the bonds' value falls
to $B$; to first order $\Delta y^* \approx (1-\ell)/D$. The second-order loss
$D\,\Delta y - \frac12\,\mathcal C\,(\Delta y)^2$ never exceeds $D^2/(2\mathcal C)$, so
when $1 - \ell > D^2/(2\mathcal C)$ the quadratic has no root and only a full
repricing gives $\Delta y^*$. After a rise leaving the bonds worth $G'$, the
fund must sell $2B - G'$ of them to restore a repo share of one half without
new capital.
\end{proposition}
\begin{proof}
The cushion is $G' - B$. To first order $G'/G = 1 - D\Delta y$; setting $G' = B$
gives the first claim. The quadratic $D x - \frac12 \mathcal C x^2$ is maximal at
$x = D/\mathcal C$ with value $D^2/(2\mathcal C)$. After selling $S$ and repaying
$S$ of repo, the share is $(B - S)/(G' - S) = \frac12$, so $S = 2B - G'$.
\end{proof}

For the long \omterm{def:m2:european-and-japanese-government-bonds:ldi}{gilt} of \cref{fig:m2:european-and-japanese-government-bonds:ldi},
$D = 21.8$ and $\mathcal C = 589$: the quadratic's largest loss is 40.5\%, less
than the fund's 50\% cushion, so the second-order formula would say the cushion
can never be exhausted. It is exhausted at 351 basis points. Taylor expansions
are for small moves; stress tests are not about small moves.

\begin{omfigure}
\begin{tikzpicture}
\begin{axis}[width=11cm,height=4.6cm,
  xlabel={rise in gilt yields (bp)},ylabel={cushion left (\%)},
  tick label style={font=\small},label style={font=\small},
  xmin=0,xmax=400,ymin=-25,ymax=105,grid=major,grid style={gray!25},
  table/col sep=comma]
\addplot[omDef,very thick] table[x=shock,y=cushion]{figdata/markets-2/07-european-and-japanese-government-bonds/ldi.csv};
\draw[omThm,thick,dashed] (axis cs:160,-25) -- (axis cs:160,105);
\draw[omProp,thick,dashed] (axis cs:250,-25) -- (axis cs:250,105);
\node[font=\footnotesize,anchor=south east,text=omThm,align=right] at (axis cs:156,-22) {September 2022:\\+160};
\node[font=\footnotesize,anchor=south west,text=omProp] at (axis cs:253,55) {standard: 250};
\end{axis}
\end{tikzpicture}

\omcaption{The cushion of an illustrative LDI fund, two-times levered in repo
and holding a thirty-year \omterm{def:m2:european-and-japanese-government-bonds:ldi}{gilt}, as \omterm{def:m2:european-and-japanese-government-bonds:ldi}{gilt} yields rise from 3.5\%. It falls to 43\%
of its starting value after 160 basis points, the size of the September 2022
move, and to zero after 351. Data: the chapter's weekend
problem.}\label{fig:m2:european-and-japanese-government-bonds:ldi}
\end{omfigure}

\begin{dated}{2026-09}{The Bank of England's gilt holdings}\label{dat:m2:european-and-japanese-government-bonds:uk}
On 17 September 2026 the Bank paused its auctions selling \omterm{def:m2:european-and-japanese-government-bonds:ldi}{gilts} from the Asset
Purchase Facility while it reviews how to continue: \omterm{def:m2:european-and-japanese-government-bonds:ldi}{gilts} maturing before 2035
(GBP~222 billion) are to be held to maturity, the longest (GBP~120 billion) held to
maturity to back the note issue, and those maturing between 2035 and 2049
(GBP~146 billion at purchase) possibly sold to the government at GBP~20 billion a
year, subject to a review before April 2027.
\end{dated}

\section{Tutorial: a spread monitor}\label{tut:m2:european-and-japanese-government-bonds:monitor}

\begin{tutorial}
\textbf{Goal.} Build euro-area spreads from monthly yields, decompose the 2011
widening, and flag unusual moves. \textbf{End state:}
\cref{fig:m2:european-and-japanese-government-bonds:spreads} and the numbers of
\cref{ex:m2:european-and-japanese-government-bonds:2011}.
\begin{tutsteps}
\item \textbf{Spreads and their decomposition.}
\omcode{code/firm/sovspread/firm_sovspread.py}{10}{31}{Spreads to a benchmark, and a change in spread split into the two yields' moves.}\label{lst:m2:european-and-japanese-government-bonds:spreads}
\item \textbf{A z-score alert} over a rolling window.
\omcode{code/firm/sovspread/firm_sovspread.py}{34}{42}{Rolling z-score of a spread.}\label{lst:m2:european-and-japanese-government-bonds:z}
\item \textbf{Run} \texttt{sovereign\_demo.widening\_2011()}: 326, 224 and 102
basis points; \texttt{fig\_sovereign.py} writes the three charts.
\item \textbf{The LDI fund.}
\omcode{code/markets-2/07-european-and-japanese-government-bonds/python/sovereign_demo.py}{51}{59}{An LDI fund's gilts, repo and cushion after a rise in yields.}\label{lst:m2:european-and-japanese-government-bonds:fund}
\end{tutsteps}
\textbf{What to change next.} Run the z-score alert on Italy's spread with a
24-month window and list the months it flags; then give the LDI fund three-times
leverage and find the yield rise that exhausts it.
\end{tutorial}

\section{Build: the sovereign spread monitor}\label{bld:m2:european-and-japanese-government-bonds:monitor}

\begin{build}
\bfield{Purpose} The miniature firm's rates desk trades euro-area spreads and
needs a live view of them, of what moved each one, and of which moves are
unusual.
\bfield{Interface} \texttt{spreads(yields, benchmark)};
\texttt{decompose(issuer, benchmark, start, end)} returning
\texttt{Decomposition(change\_bp, issuer\_bp, benchmark\_bp)};
\texttt{zscore(series, window)}; \texttt{alerts(series, window, threshold)}.
\bfield{Rules} Yields in percent, spreads in basis points; spreads only on
dates both issuers have; the z-score of a value uses the window
\emph{before} it.
\bfield{Acceptance tests} \texttt{code/firm/sovspread/tests/}: spreads on
common dates; the decomposition adds up and gives the benchmark's fall a
widening sign; a jump triggers the alert and noise does not.
\bfield{Stretch} Daily data from the issuers' own benchmark curves; spreads
at matched maturities with the bond library
(\cref{bld:m2:government-bonds:bond}); a spread in asset-swap terms
(\cref{ch:m2:corporate-bonds}).
\end{build}

\begin{omsources}
\item Bank of England, letter from Sir Jon Cunliffe to the Treasury Committee on
LDI, 5 October 2022; \emph{Quarterly Bulletin} 2023, ``Financial stability buy/sell
tools: a gilt market case study''; staff paper on LDI minimum resilience, 2023;
APF market notice, 17 September 2026.
\item European Central Bank, press releases on outright monetary transactions
(6 September 2012) and the transmission protection instrument (21 July 2022).
\item Bank of Japan, ``New Framework for Strengthening Monetary Easing'', 21
September 2016; ``Changes in the Monetary Policy Framework'', 19 March 2024.
\item OECD, long-term interest rates (via FRED); MTS, company website.
\end{omsources}

\section{Exercises}

\begin{exercise}[$\star$]\label{exo:m2:european-and-japanese-government-bonds:1}
With the yields of \cref{dat:m2:european-and-japanese-government-bonds:yields},
give the spreads of Spain, Italy and France to Germany in basis points.
\end{exercise}

\begin{exercise}[$\star$]\label{exo:m2:european-and-japanese-government-bonds:2}
Italy's yield rises 30 basis points and Germany's falls 10 in a day. Give the
change in the spread and its decomposition.
\end{exercise}

\begin{exercise}[$\star$]\label{exo:m2:european-and-japanese-government-bonds:3}
An LDI fund holds GBP~1 billion of \omterm{def:m2:european-and-japanese-government-bonds:ldi}{gilts}, half financed in repo. What is its
leverage? What is its cushion after a 20\% fall in \omterm{def:m2:european-and-japanese-government-bonds:ldi}{gilt} prices, and its new
leverage?
\end{exercise}

\begin{exercise}[$\star\star$]\label{exo:m2:european-and-japanese-government-bonds:4}
With \cref{prop:m2:european-and-japanese-government-bonds:cushion}, estimate the
yield rise that exhausts the cushion of a fund with duration 21.8, convexity 589
and half its \omterm{def:m2:european-and-japanese-government-bonds:ldi}{gilts} in repo, to first order. What does the second-order formula
say, and why?
\end{exercise}

\begin{exercise}[$\star\star$]\label{exo:m2:european-and-japanese-government-bonds:5}
Why did the Bank of England buy only long \omterm{def:m2:european-and-japanese-government-bonds:ldi}{gilts} in 2022, and only for a
fortnight? Contrast with \omterm{def:m2:central-banks-and-the-short-rate:qe}{quantitative easing}.
\end{exercise}

\begin{exercise}[$\star\star$]\label{exo:m2:european-and-japanese-government-bonds:6}
Under \omterm{def:m2:european-and-japanese-government-bonds:ycc}{yield-curve control}, the market sells ten-year bonds at the target
yield. Who decides how many bonds the central bank buys? What happens to its
balance sheet if the market expects the target to be abandoned?
\end{exercise}

\begin{exercise}[$\star\star\star$]\label{exo:m2:european-and-japanese-government-bonds:7}
\emph{Coding.} With \texttt{sovereign\_demo}, compute the sales the LDI fund
must make to restore half-leverage after rises of 100 and 160 basis points,
and the cushion (as a share of the holding) needed to survive 250.
\end{exercise}

\begin{exercise}[$\star\star\star$]\label{exo:m2:european-and-japanese-government-bonds:8}
\emph{Find the flaw.} ``The LDI funds lost money in 2022 because rising yields
made the pension schemes poorer.'' Correct the statement using the Bank's own
account of the schemes' balance sheets.
\end{exercise}

\section{Problem: The Collateral Spiral}

\begin{problem}[{Weekend problem --- how far an LDI fund can fall}]\label{pb:m2:european-and-japanese-government-bonds:1}
On 22 September 2022 an LDI fund holds GBP~1 billion (market value) of a
thirty-year \omterm{def:m2:european-and-japanese-government-bonds:ldi}{gilt} with a 1.5\% coupon, maturing 22 July 2052, at a yield of
3.50\%; half is financed in repo. The fund cannot raise new capital in less
than a week. The \omterm{def:m2:european-and-japanese-government-bonds:ldi}{gilt} and the fund are illustrative.

\medskip\noindent\textbf{Part I --- The fund.}
\begin{enumerate}
\item Give the \omterm{def:m2:european-and-japanese-government-bonds:ldi}{gilt}'s dirty price, the fund's repo and its cushion.
\item Give the \omterm{def:m2:european-and-japanese-government-bonds:ldi}{gilt}'s \omterm{def:m2:government-bonds:duration}{modified duration} and the fund's DV01 per basis point.
\item What is the fund's leverage, and its DV01 per pound of capital?
\item How many basis points of rise cost 1\% of the cushion, to first order?
\item Why would a pension scheme choose such a fund rather than hold \omterm{def:m2:european-and-japanese-government-bonds:ldi}{gilts}
outright?
\end{enumerate}

\medskip\noindent\textbf{Part II --- The shock.}
\begin{enumerate}[resume]
\item Yields rise 100 basis points. Give the \omterm{def:m2:european-and-japanese-government-bonds:ldi}{gilts}' value, the cushion, and the
leverage.
\item The same for 160 basis points, the size of the September move.
\item How much must the fund sell after 100 basis points to restore
half-leverage without new capital? After 160?
\item Compare the aggregate sales that managers reported with the market's daily
volume. What happens to the price of what they sell?
\item Why did the collateral calls come on swaps as well as repo?
\end{enumerate}

\medskip\noindent\textbf{Part III --- The limit.}
\begin{enumerate}[resume]
\item Give the rise that exhausts the cushion.
\item What happens to the repo lender when it is exhausted?
\item Give the cushion the fund needs to survive 250 basis points without
selling.
\item Earlier stress tests used 100 basis points. What cushion did that imply?
\item Why is a percentage cushion not enough to state resilience, without the
duration of what is held?
\end{enumerate}

\medskip\noindent\textbf{Part IV --- Judgement.}
\begin{enumerate}[resume]
\item Why did a central bank with an inflation problem buy long bonds?
\item Why were the purchases temporary and small against the market?
\item Why was the pension scheme itself, as opposed to the fund, not insolvent?
\item State the \emph{named result}: the yield rise that exhausts the cushion,
and the cushion after the September move.
\item In one sentence: what turned a rise in yields into a crisis?
\end{enumerate}
\end{problem}

\section{Interview questions}

\begin{interviewq}[$\star$ \iqroles{trader, researcher}]\label{iq:m2:european-and-japanese-government-bonds:1}
Why do Italian and German ten-year bonds, in the same currency, trade at
different yields? What is in the spread?
\end{interviewq}

\begin{interviewq}[$\star$ \iqroles{trader, bank}]\label{iq:m2:european-and-japanese-government-bonds:2}
What happened in the \omterm{def:m2:european-and-japanese-government-bonds:ldi}{gilt} market in September 2022?
\end{interviewq}

\begin{interviewq}[$\star\star$ \iqroles{trader, researcher}]\label{iq:m2:european-and-japanese-government-bonds:3}
How does \omterm{def:m2:european-and-japanese-government-bonds:ycc}{yield-curve control} work, and what are the risks of exiting it?
\end{interviewq}

\begin{interviewq}[$\star\star$ \iqroles{researcher, trader}]\label{iq:m2:european-and-japanese-government-bonds:4}
A spread widened by 100 basis points. How do you tell how much was the issuer
and how much was the benchmark, and why does it matter for a hedge?
\end{interviewq}

\begin{interviewq}[$\star\star$ \iqroles{bank, researcher}]\label{iq:m2:european-and-japanese-government-bonds:5}
What is the difference between the ECB's OMT and TPI?
\end{interviewq}

\begin{interviewq}[$\star\star\star$ \iqroles{researcher, bank}]\label{iq:m2:european-and-japanese-government-bonds:6}
You are asked to stress-test a leveraged bond fund. How do you choose the size
of the shock, and what else must the test capture?
\end{interviewq}
